% Sect. 5 Benchmark (outline brief 6.5; mirrors RQ1 Sect. 5: an opening paragraph, then the run-in heads
% "Development and test cases.", "Inputs." (new), "Reference deltas." (RQ1: "Reference sets.") and "Matching.").
% Test cases: nothing beyond what the RQ1 paper reports (names, RQ1 Table 1 values). No base release and no
% reference-profile version is chosen here (Q3). The reference deltas are not built yet: their counts are \res slots.
% Source shorthand in the % src: comments: outline-final.md Sect. 0.1 (B = backend/profile_extensioner/).
\section{Benchmark}\label{sec:benchmark}

The benchmark adds to the five cases of the companion study \cite{anon2027rq1} what each published extension adds
to its base: a \emph{reference delta}, the published profile minus its base (Table~\ref{tab:cases}). % src: RQ1 L358, L381-388 (five cases); PLAN:9-11
Appendix~\ref{app:benchmark} has the cases in full and, in Fig.~\ref{fig:reference}, the path from a published
profile to Recovery. Each case runs against one pinned base release, one of the nine the tool holds
(Sect.~\ref{sec:impl}), and the \ac{DCAT} version it builds on, because a reference delta exists only relative to
one release. % src: FACTS5:15-23 (nine bases); DCS §3
Against \ac{DCAT-AP}~1.1, built on DCAT~2014, which has no data services, the data service that \ac{DCAT-AP}~2.0.0
profiles is an addition. % src: @5ffd56e B/data/base-profiles/dcat-ap-2.0.0/profile.yaml:25-30; B/data/base-profiles/dcat-ap-1.1/profile.yaml:21-25; B/data/dcat/history.json:30-38 (dcat:DataService introduced in dcat-2)
\todo{base and reference version per case (Q3)} Every reference delta is partial (shape constructs the comparison
does not handle, profiles that publish no shapes).

\begin{table}[tb]
\caption{The benchmark cases, test cases above development cases. Base in~\cite{anon2027rq1}: the companion % src: RQ1 L377, L381-388
study's \acs{DCAT-AP} release. Base: the release the reference delta is computed against, with its \acs{DCAT}
version. Reference: the published profile and version. Atoms: the companion's scored content atoms, the input.
Elements: reference elements; aligned ($N$): those at least one atom asks for, the denominator.}
\label{tab:cases}
\centering\scriptsize
% Atoms and the companion's bases: RQ1 Table 1 (RQ1 L381-388). Base and Reference: the author decides (Q3).
% Elements and Aligned (N): written by the RQ2 evaluator (\res prints \tbd until then).
% Bases, references and alignments are fixed before any evaluation run (stated in the text, Development and test cases).
\begin{tabular}{@{}l l@{\hspace{16pt}}l l r r r@{}}
\toprule
Case & Base in~\cite{anon2027rq1} & Base & Reference (version) & Atoms & Elements & Aligned ($N$) \\
\midrule
mobilityDCAT-AP & 2.0.1 & \todo{Q3} & \todo{Q3} & 39 & \res{mobility}{all}{elements} & \res{mobility}{all}{N} \\ % src: RQ1 L382
GeoDCAT-AP & 1.1 & \todo{Q3} & \todo{Q3} & 37 & \res{geo}{all}{elements} & \res{geo}{all}{N} \\ % src: RQ1 L383
HealthDCAT-AP & 3.0.0 & \todo{Q3} & \todo{Q3} & 89 & \res{health}{all}{elements} & \res{health}{all}{N} \\ % src: RQ1 L384
\emph{Test cases} & & & & 165 & \res{test}{all}{elements} & \res{test}{all}{N} \\ % src: RQ1 L385 (39+37+89)
\midrule
StatDCAT-AP & 1.1 & \todo{Q3} & \todo{Q3} & 14 & \res{stat}{all}{elements} & \res{stat}{all}{N} \\ % src: RQ1 L386
DCAT-AP-JRC & 1.1 & \todo{Q3} & \todo{Q3} & 17 & \res{jrc}{all}{elements} & \res{jrc}{all}{N} \\ % src: RQ1 L387
\midrule
All five & & & & 196 & \res{five}{all}{elements} & \res{five}{all}{N} \\ % src: RQ1 L388 (165+14+17)
\bottomrule
\end{tabular}
\end{table}

\paragraph{Development and test cases.} DCAT-AP-JRC and StatDCAT-AP, the companion study's \emph{development
cases}, are those the pipeline was built on, reported apart (Sect.~\ref{sec:results}); GeoDCAT-AP, HealthDCAT-AP
and mobilityDCAT-AP are the \emph{test cases}. No test case was run, no reference delta
was built for one, and no test-case atom or output was seen. The design drew on the publication practice of
published extensions, including newer releases of the test-case profiles, for 18 of the 25 profile-practice steps;
these are cited through the guidance they restate. Table~\ref{tab:steps} in the appendix marks them. % src: LIT:74-98 (25 G3 rows; 18 name a test-case key: S01, S02, S04, S06, S07, S08, S09, S10, S11, S13, S14, S16, S17, S19, S20, S23, S24, S25); outline-final 7.3; ADD §H (binding wording)
\todo{test-case statement: confirm (Q18)} The pipeline, the reference deltas and their alignments are frozen
before any case is scored, and all five cases are evaluated alike. % src: RQ1 L391-392; PLAN:19-20 (design)
\todo{freeze: commit and tag} No part of the pipeline was written for one case.
\verify{1.1 base, importer, recommended-minimum normalization added during development}
% The verify: the 1.1 base, the reference importer and the normalization of a recommended minimum were added while
% the development cases ran.
% The plan freezes reference deltas, atoms, alignments and counterfactual variants as one git tag (PLAN:19-20).

\paragraph{Inputs.} The input of each case is its \emph{reference atoms}: every content atom the companion
benchmark scores for it, not only those its corpus states, because the pipeline reads no corpus
(Table~\ref{tab:cases}). \todo{test-case atom counts: print? (Q11)} An importer keeps each atom's identifier,
statement, subject, modality and source reference and drops the evaluator fields (relation to the base, anchors,
evidence labels); these atoms isolate the modeling step.
% src: @5ffd56e scripts/import_rq1_reference.py:1-16, 47-59; RQ1 L403-410 (reference atoms: the plan's A-gold)
The companion's reference sets record the source's kind of statement, not a modality. The importer maps explicit
requirements and mandatory items to must, recommendations to should, optional items to may and the rest to
unspecified, which licenses at most recommended (Sect.~\ref{sec:problem}).
% src: @5ffd56e scripts/import_rq1_reference.py:12-14, 26-32, 56; @5ffd56e B/contracts.py:713-715; PILOT:4-5
\verify{share of unspecified modalities per case} \todo{modality from the reference's obligation? (Q20)} A second
input, the \emph{extracted atoms}, converted from the needs a run of the companion workflow retained (it accepted
none), shows how extraction errors carry into the profile. % src: RQ1 L474-475; PLAN:24-25 (extracted atoms: the plan's A-rq1)
\todo{extracted atoms: run, conversion, in or out (Q21)}

\paragraph{Reference deltas.} Each reference delta compares the published profile's shapes and specified
obligations with its base's, per class and property, both merged as in Sect.~\ref{sec:problem}. An element is
\emph{added} where the base states nothing, \emph{tightened} where the profile is stricter on a facet and
\emph{loosened} where it is weaker. Loosened elements are reported apart; unchanged ones are dropped, because an
inherited constraint earns no credit.
% src: PLAN:13-15; PE:1-28, :127-128 (DeltaStatus); the delta spike's docstring (Q23)
Class additions are kept, usage notes are not scored, and the constraints on one class and path fold into one
element. Unhandled constructs are listed, never scored; a profile without shapes is extracted by hand from its
tables, and a second person checks a sample. \todo{reference deltas: source and construction (Q23)} Each element
is labeled with its kind, facets and the atoms that ask for it. Elements no atom asks for stay outside the
denominator, the \emph{aligned reference elements}~$N$ (Sect.~\ref{sec:problem}), fixed before any change is
matched. \todo{alignment and hand extraction: who, agreement (Q9)}

\paragraph{Matching.} Generated and reference deltas share one contract, the ProfileDelta
(Sect.~\ref{sec:problem}). A generated and a reference element match \emph{strictly} when class and path
\acp{IRI} agree, and in the \emph{relaxed} sense when one path is an RDFS sub-property of the other and an
adjudicator confirms the pair. A minted element matches \emph{by meaning}, on its class and definition,
never by \ac{IRI}. % src: PLAN:67-68; EXP:24-26; PE:3-28
A reference element is credited once, by its best match in that order. Matched pairs are compared on obligation,
minimum, maximum, range and value list, not on severity, which follows the obligation. % src: PLAN:69; KCU §8 item 12; ISS K12 (@5ffd56e B/compile/shacl.py:76; B/engine/propose.py:295)
\emph{Recovery} counts the aligned reference elements with any match, and \emph{Strict} those with a strict match.
Recovery is thus a recall over the aligned reference elements, not a share of correct changes or of the published
profile. Recovery leads, with Strict beside it in every figure: relaxed matches and matches by meaning rest on a
judgment but still point a reviewer to the element. Unmatched included changes are adjudicated as an alternative
modeling of an aligned atom, as beyond the reference, or as wrong. Adjudicators see the distinct elements of a
case's runs pooled under random identifiers, without run or condition. % src: PLAN:65-69 (design)
\todo{adjudicator and agreement (Q9)} \todo{evaluator: not built}
